\origpage{82}{ROBERTSON AND SEYMOUR}
\noindent we may choose $r_1, r_2$ both incident with some edge of $\Gamma$ not in $E(N(Z))$. But this contradicts the definition of $N(Z)$, since $r_1 \in Z$, $r_2 \notin Z$, and completes the proof. $\square$

\stmt
(5.3) \textit{If $Z \subseteq R(\Gamma)$ is small, $Z' \subseteq Z$, and $|\Omega \cap V(N(Z'))| < \omega$, then $Z'$ is small.}

\textit{Proof.} We must show that $N(Z')$ captures every $r \in Z'$. Thus, let $r \in Z'$. Then $r \in Z$, and so $N(Z)$ captures $r$; let $C$ be a circuit of $N(Z)$ with $r \subseteq \operatorname{ins}(C)$. Let $r'$ be the region of $N(Z')$ with $r \subseteq r'$. We claim that $r' \subseteq \operatorname{ins}(C)$. Suppose not. Let $X = \{s \in R(\Gamma) : s \subseteq \operatorname{ins}(C)\}$. Since $X \cap Z'$ contains some but not every region of $\Gamma$ included in $r'$, by (5.1) there are regions $r_1, r_2$ of $\Gamma$ with $r_1, r_2 \subseteq r'$, and with $r_1 \in X \cap Z'$, $r_2 \notin X \cap Z'$, both incident with some edge $e$ of $\Gamma$ not in $N(Z')$. Since $e \notin E(N(Z'))$ and $r_1 \in Z'$ it follows that $r_2 \in Z'$, and so $r_2 \notin X$. Since $r_1 \in X$ we deduce that $e \in E(C) \subseteq E(N(Z))$. But this is impossible since $r_1, r_2 \in Z' \subseteq Z$. This proves our claim that $r' \subseteq \operatorname{ins}(C)$, and so $\Gamma$ captures $r'$. By (4.5), $N(Z')$ captures $r'$ and hence $r$, as required. $\square$

The main and most difficult result of this section is the following. (We say $(Z_1, Z_2, Z_3)$ is a partition of $Z$ if $Z_1, Z_2, Z_3$ are disjoint sets whose union is $Z$.)

\stmt
(5.4) \textit{If $(Z_1, Z_2, Z_3)$ is a partition of $R(\Gamma)$, then one of $Z_1, Z_2, Z_3$ is not small.}

\textit{Proof.} We denote $N(Z_1) \cup N(Z_2) \cup N(Z_3)$ by $N(Z_1, Z_2, Z_3)$. If possible, let $(Z_1, Z_2, Z_3)$ be a counterexample chosen with $N(Z_1, Z_2, Z_3)$ minimal. Then $Z_1$, $Z_2$, $Z_3$ are small. Therefore for $i = 1, 2, 3$ we have $|\Omega \cap V(C)| < \omega$ and consequently $|V(C)| < 2\theta$ for every circuit $C$ in $N(Z_i)$, from which it follows that $N(Z_i)$ is confined and $\operatorname{ins}(C)$ is defined for every circuit $C$ in $N(Z_i)$. Moreover $Z_1 \ne \emptyset$ because otherwise a contradiction would be obtained by taking $Z = Z_2$ in (5.2), and similarly $Z_2, Z_3$ are non-empty.

A $Z_1$-ring is a circuit $C$ of $N(Z_1)$ chosen with $\operatorname{ins}(C)$ maximal. If $C_1$ is a $Z_1$-ring and $C$ is a circuit of $N(Z_1)$ with $|V(C) \cap \operatorname{ins}(C_1)| \ge 2$, then $U(C) \subseteq \operatorname{ins}(C_1)$ by (4.3)(i). Consequently, if $C_1$, $C_2$ are distinct $Z_1$-rings and $C$ is a circuit of $N(Z_1)$ with $|V(C \cap C_i)| \ge 2$ ($i = 1, 2$), then $|V(C_2) \cap \operatorname{ins}(C_1)| \ge 2$ since $U(C) \subseteq \operatorname{ins}(C_1)$; and so $U(C_2) \subseteq \operatorname{ins}(C_1)$, and consequently $\operatorname{ins}(C_2) \subseteq \operatorname{ins}(C_1)$ by the first slope axiom. But similarly $\operatorname{ins}(C_1) \subseteq \operatorname{ins}(C_2)$, and so equality holds, which is impossible since $C_1 \ne C_2$. Thus there do not exist such $C_1, C_2, C$. In particular, any two distinct $Z_1$-rings intersect in at most one vertex and hence have no common edge, and also if $e, f$ are edges of distinct $Z_1$-rings then no circuit of $N(Z_1)$
